\documentclass[11pt,a4paper]{article}

% =============================================================================
% Packages
% =============================================================================
\usepackage[utf8]{inputenc}
\usepackage[T1]{fontenc}
\usepackage[english]{babel}
\usepackage{amsmath,amssymb,amsthm}
\usepackage{geometry}
\geometry{margin=2.6cm}
\usepackage{booktabs}
\usepackage{longtable}
\usepackage{array}
\usepackage{enumitem}
\usepackage{tikz}
\usepackage{float}
\usepackage{hyperref}
\hypersetup{
  colorlinks=true,
  linkcolor=blue!60!black,
  urlcolor=blue!60!black,
  citecolor=blue!60!black,
  pdftitle={The 1D Hydraulic Calculation Process of the LocalSWMM Engine}
}
\usepackage{microtype}
\usepackage{xcolor}
\usepackage{siunitx}

% =============================================================================
% Environment & macros
% =============================================================================
\newtheorem{definition}{Definition}[section]
\newtheorem{remark}{Remark}[section]

\newcommand{\st}{\mathrm{s}}
\newcommand{\ft}{\mathrm{ft}}
\newcommand{\cfs}{\mathrm{ft^3/s}}
\newcommand{\sfft}{\mathrm{ft^2}}
\newcommand{\R}{\mathbb{R}}
\renewcommand{\arraystretch}{1.25}

% Shorthand for the engine paths
\newcommand{\src}[1]{\texttt{src/engine/#1}}

\title{\textbf{The 1D Hydraulic Calculation Process}\\[0.6em]
\large of the LocalSWMM project (HydroCouple OpenSWMM engine)}
\author{Technical Reference Document}
\date{August 2026}

\begin{document}

\maketitle

\begin{abstract}
This document provides a complete, equation-level description of the
one-dimensional (1D) hydraulic calculation process used by the LocalSWMM web
application. All hydraulic computations are executed by the HydroCouple
\textbf{OpenSWMM} engine, which is compiled to WebAssembly and runs inside the
browser. The engine solves the de~Saint-Venant equations on a node--link
network. Four routing formulations are available: \emph{steady flow},
\emph{kinematic wave}, \emph{dynamic wave} (the default), and an explicit
\emph{finite-volume} scheme. This document explains the conceptual model, the
governing equations, the finite-difference schemes, the iterative solution
procedure, the node-continuity and surcharge algorithms, the computation of
non-conduit hydraulic structures, the cross-section geometry functions,
boundary conditions, adaptive time stepping, and the mass-balance accounting.
Every equation is transcribed from the engine source code (see the
References for the file list) and cross-referenced to the engine's Reference
Manuals where applicable.
\end{abstract}

\tableofcontents
\newpage

% =============================================================================
% 1. INTRODUCTION
% =============================================================================
\section{Introduction}

\subsection{Scope}

This document covers the \emph{one-dimensional} hydraulic (flow-routing)
calculations of the LocalSWMM project. LocalSWMM is a browser-based web
application for stormwater and wastewater network modeling; it embeds the
HydroCouple OpenSWMM engine (C$++$) compiled to WebAssembly. The 1D hydraulic
component solves unsteady, gradually varied, one-dimensional flow through a
network of \emph{nodes} (junctions, storage units, dividers, outfalls) and
\emph{links} (conduits, pumps, orifices, weirs, outlets). The 2D overland-flow
module and its 1D--2D coupling are out of scope except where they interact with
the 1D solver (e.g.\ through lateral inflows and the ponded-area coupling).

The engine is a \emph{distributed discrete-time simulation model}: it advances
a state vector over a sequence of time steps,
\begin{equation}
  \boldsymbol{X}_{t} = f\!\left(\boldsymbol{X}_{t-1}, \boldsymbol{I}_{t}, \boldsymbol{P}\right),
  \qquad
  \boldsymbol{Y}_{t} = g\!\left(\boldsymbol{X}_{t}, \boldsymbol{P}\right),
\end{equation}
where $\boldsymbol{X}$ is the state vector, $\boldsymbol{I}$ external inputs,
$\boldsymbol{P}$ fixed parameters, and $\boldsymbol{Y}$ the outputs.

\subsection{The node--link conceptual model}

The conveyance portion of a drainage system is a network of nodes and links.
\textbf{Nodes} are points that represent:
\begin{itemize}[nosep]
  \item \textbf{Junctions} --- points where links join, with invert elevation,
        height-to-rim, optional surcharge (pressure) depth, and optional
        ponded surface area.
  \item \textbf{Outfalls} --- terminal boundary nodes with a prescribed stage
        condition (free, normal, fixed, tidal, or time-series) and an optional
        flap gate.
  \item \textbf{Storage units} --- nodes with a surface-area--versus--depth
        relation and real storage volume; they never pressurize.
  \item \textbf{Flow dividers} --- nodes that split inflow between a
        continuation link and a diversion link (cutoff, overflow, tabular, or
        weir; active under kinematic wave, treated as junctions under dynamic
        wave).
\end{itemize}

\textbf{Links} connect nodes. The link types are:
\begin{itemize}[nosep]
  \item \textbf{Conduits} --- pipes and open channels of arbitrary
        cross-section, described by Manning's equation and a set of cross-section
        geometry tables.
  \item \textbf{Pumps} --- governed by pump curves (five curve types or an
        ideal-pump model) with on/off depth hysteresis.
  \item \textbf{Orifices} --- bottom or side openings with weir/orifice
        discharge behavior.
  \item \textbf{Weirs} --- transverse, side-flow, V-notch, or trapezoidal weirs
        with an optional transition to orifice flow under submergence.
  \item \textbf{Outlets} --- head or depth rating curves (tabular or
        functional).
\end{itemize}

\subsection{State variables and internal units}

Under flow routing the fundamental state variables are exactly three
(Table~\ref{tab:statevars}):

\begin{table}[H]
\centering
\begin{tabular}{ll}
\toprule
\textbf{Symbol} & \textbf{Description} \\
\midrule
$H$  & Hydraulic head of water at a node \\
$Q$  & Flow rate in a link \\
$A$  & Flow area in a link (inferred from $Q$ and geometry) \\
\bottomrule
\end{tabular}
\caption{The three fundamental flow-routing state variables.}
\label{tab:statevars}
\end{table}

All other quantities are derived from these three variables, the external
inputs, and the fixed input parameters. Internally the engine carries out all
calculations in \textbf{feet} (length), \textbf{seconds} (time), \textbf{cfs}
(flow), \textbf{ft$^2$} (area), and \textbf{ft$^3$} (volume). Project input
values in SI (metric) or US units are converted to these internal units at
parse time; the conversion factors are
\begin{align}
  \mathrm{Ucf}[\mathrm{LENGTH}] &=
    \begin{cases} 1.0 & \text{(US)} \\ 0.3048 & \text{(SI)} \end{cases}, \\
  \mathrm{Ucf}[\mathrm{VOLUME}] &=
    \begin{cases} 1.0 & \text{(US)} \\ 0.02832 & \text{(SI)} \end{cases},
\end{align}
(note that the SI volume factor is the legacy \emph{truncated} value, not
$0.3048^3$, deliberately preserved for floating-point parity with the legacy
SWMM engine). Flow conversion uses $\mathrm{Qcf} = \{1.0,\ 448.831,\ 0.64632,\
0.02832,\ 28.317,\ 2.4466\}$ for cfs/gpm/MGD/cms/LPS/MLD.

\subsection{The four routing formulations}

The engine selects the flow-routing method with the \texttt{FLOW\_ROUTING}
option. Table~\ref{tab:routing} summarizes the four methods.

\begin{table}[H]
\centering
\begin{tabular}{@{}l p{4.6cm} p{5.4cm}@{}}
\toprule
\textbf{Method} & \textbf{Governing equations} & \textbf{Features} \\
\midrule
Steady flow & continuity + Manning normal depth & pass-through, no storage/backwater \\
Kinematic wave & continuity + uniform-flow rating & no backwater, reversal, or surcharge \\
Dynamic wave (default) & full St.\ Venant & backwater, losses, reversal, surcharge \\
Finite volume (explicit) & St.\ Venant, Godunov/HLLC & transcritical flow, native dry/wet \\
\bottomrule
\end{tabular}
\caption{The four 1D routing formulations of the engine.}
\label{tab:routing}
\end{table}

The dynamic wave is the default and is treated in the greatest depth here. Each
routing method is invoked from a common \texttt{Router} orchestrator.

% =============================================================================
% 2. MODEL COMPONENTS
% =============================================================================
\section{Hydraulic model components}

\subsection{Node data}

Each node carries static geometric properties and dynamic state
(\texttt{NodeData.hpp}):
\begin{itemize}[nosep]
  \item $z_{\mathrm{inv}}$ --- invert elevation (ft).
  \item $d_{\mathrm{full}}$ --- full depth (depth to the rim / top of storage).
  \item $d_{\mathrm{sur}}$ --- maximum allowed surcharge depth above full
        depth; the flooding limit is $z_{\mathrm{inv}} + d_{\mathrm{full}} +
        d_{\mathrm{sur}}$.
  \item $A_{\mathrm{pond}}$ --- ponded surface area used when the node floods
        above its rim and ponding is enabled.
  \item $z_{\mathrm{crown}}$ --- elevation of the crown of the highest
        connecting conduit; the surcharge threshold is $y_{\mathrm{crown}} =
        z_{\mathrm{crown}} - z_{\mathrm{inv}}$.
  \item $\mathrm{deg}$ --- number of connecting links; if negative the node is
        an upstream terminal node.
  \item $V_{\mathrm{full}}$ --- volume at full depth.
\end{itemize}

Dynamic state (updated each routing step):
\begin{itemize}[nosep]
  \item $d$ --- water depth above the invert (ft);
  \item $H = z_{\mathrm{inv}} + d$ --- hydraulic head (ft);
  \item $V$ --- stored volume (ft$^3$);
  \item $Q_{\mathrm{lat}}$ --- total lateral inflow (ft$^3$/s);
  \item $Q_{\mathrm{in}}$, $Q_{\mathrm{out}}$ --- total inflow and outflow;
  \item $Q_{\mathrm{ovfl}}$ --- overflow/flooding rate;
  \item $Q_{\mathrm{loss}}$ --- evaporation + seepage/exfiltration loss rate;
  \item $\Delta Q_{\mathrm{net}}^{t-1}$ --- net inflow $Q_{\mathrm{in}} -
        Q_{\mathrm{out}}$ of the previous step (used for trapezoidal
        averaging);
  \item previous-step values $d^{t-1}$, $V^{t-1}$, $Q_{\mathrm{lat}}^{t-1}$.
\end{itemize}

\subsubsection{Node surface area and volume functions}

The node geometry functions are implemented in \src{hydraulics/Node.cpp}.

\textbf{Junction/outfall/divider volume}: with $\mathrm{MIN\_SURFAREA} =
12.566\ \sfft$ ($\approx 4\pi$, the area of a 4-ft manhole),
\begin{equation}
  V(d) = V_{\mathrm{full}}\,\frac{d}{d_{\mathrm{full}}},
\end{equation}
where $V_{\mathrm{full}} = \mathrm{MIN\_SURFAREA}\cdot d_{\mathrm{full}}$
unless overridden.

\textbf{Storage volume}:
\begin{itemize}[nosep]
  \item Tabular: $V = \mathrm{table}(d\cdot \mathrm{Ucf}[L])/\mathrm{Ucf}[V]$
        with table interpolation;
  \item Geometric shapes (cylindrical, conical, paraboloid, pyramidal) use the
        quadratic surface-area relation $A(d) = c + a\,d + b\,d^{2}$, integrated
        to the cubic
        \begin{equation}
          V(d) = d\left(c + d\!\left(\tfrac{a}{2} + d\,\tfrac{b}{3}\right)\right);
        \end{equation}
  \item Functional storage: power law $A(d) = c + a\,d^{b}$ integrated
        analytically.
\end{itemize}

\textbf{Surface area} $A_{\mathrm{surf}}(d)$: non-storage nodes return 0 (the
minimum-area floor is applied later in the dynamic-wave node update); storage
nodes return the table value (extrapolating linearly) or the analytic
$A(d) = c + a\,d + b\,d^{2}$ / $A(d) = c + a\,d^{b}$.

\textbf{Ponded area}:
\begin{equation}
  A_{\mathrm{pond}}(d) =
  \begin{cases}
    A_{\mathrm{surf}}(d) & d \le d_{\mathrm{full}} \text{ or } A_{\mathrm{pond}} = 0,\\
    A_{\mathrm{pond}} & d > d_{\mathrm{full}} \text{ (flooded)}.
  \end{cases}
\end{equation}

\textbf{Depth from volume} (the inverse of the volume function, used by the
kinematic-wave storage update and by reporting): junction nodes use
$d = V/\mathrm{MIN\_SURFAREA}$; storage nodes invert the tabular/geometric/
functional relations, using Newton--Raphson with bisection fallback where no
closed form exists.

\textbf{Maximum outflow} (used by the kinematic-wave and steady-flow link
inflow gathering):
\begin{equation}
  Q_{\max} = Q_{\mathrm{in}} + \frac{V_{\mathrm{old}}}{dt},
  \qquad Q_{\mathrm{in}} = \min\!\left(Q_{\mathrm{in}}, Q_{\max}\right).
\end{equation}

\subsection{Link data}

Each link carries static properties and dynamic state (\texttt{LinkData.hpp}).
For conduits the per-link (per-barrel) conveyance parameters are
(\texttt{Link.cpp}, \texttt{computeConveyance}):
\begin{align}
  \beta &= \frac{\mathrm{PHI}\,\sqrt{|S_0|}}{n}, \qquad \mathrm{PHI} = 1.486
    \quad\text{(Manning US factor)}, \\
  \mathrm{rough\_factor} &= g\left(\frac{n}{\mathrm{PHI}}\right)^{2},\\
  Q_{\mathrm{full}} &= S_{\mathrm{full}}\cdot\beta,
  \qquad Q_{\max} = S_{\max}\cdot\beta,
\end{align}
where $n$ is the Manning roughness, $S_0$ the conduit slope, $S(a) = A\,R^{2/3}$
the \emph{section factor} (see Section~\ref{sec:xsect}), $S_{\mathrm{full}} =
S(A_{\mathrm{full}})$ and $S_{\max} = \max_a S(a)$.

Dynamic state per link: $Q$ (current flow, $+$ = node1$\to$node2), $d$
(midpoint depth), $V$ (volume), $\mathrm{Fr}$ (Froude number), $\mathrm{flow\_class}$
(DRY, UP\_DRY, DN\_DRY, SUBCRITICAL, SUPERCRITICAL, UP\_CRITICAL, DN\_CRITICAL),
and previous-step values.

\subsection{Conduit slope and length}

For a conduit connecting node $i$ and node $j$ with end offsets $z_1, z_2$,
the end invert elevations of the conduit are $Z_1 = z_{\mathrm{inv},1} +
z_1$, $Z_2 = z_{\mathrm{inv},2} + z_2$. The conduit slope is
\begin{equation}
  S_0 = \frac{\Delta y}{\Delta x},
  \qquad \Delta y = Z_1 - Z_2, \quad
  \Delta x = \sqrt{L^{2} - \Delta y^{2}},
\end{equation}
with a minimum imposed drop of $\lvert \Delta y \rvert \ge 0.001\ \ft$
(overrideable).

\subsection{Conduit lengthening (Courant stability)}

If a positive \texttt{LENGTHENING\_STEP} is supplied, short conduits are
artificially lengthened so that the Courant condition for the user-supplied
time step is met (\texttt{Routing.cpp:135--195}):
\begin{equation}
  t_{\mathrm{step}} = \min(\mathrm{routing\_step},\ \mathrm{lengthening\_step}),
  \qquad
  v_{\mathrm{full}} = \frac{\mathrm{PHI}}{n}\,\frac{S_{\mathrm{full}}}
    {\sqrt{|S_0|}\,A_{\mathrm{full}}},
\end{equation}
\begin{equation}
  \mathrm{ratio} = \frac{\left(\sqrt{g\,y_{\mathrm{full}}} + v_{\mathrm{full}}\right)
    t_{\mathrm{step}}}{L},
  \qquad
  L' = \max(L,\ \mathrm{ratio}\cdot L),
\end{equation}
For open channels $y_{\mathrm{full}}$ is replaced by the hydraulic depth
$A_{\mathrm{full}}/W_{\max}$. If $L' > L$, the slope and roughness are adjusted
to preserve equal head loss at any flow,
\begin{equation}
  S_0' = \frac{|S_0|}{L'/L}, \qquad
  n' = \frac{n}{\sqrt{L'/L}},
\end{equation}
and $\beta$, $\mathrm{rough\_factor}$, $Q_{\mathrm{full}}$, $Q_{\max}$ are
recomputed from the adjusted values. The lengthened length $L'$
(\texttt{mod\_length}) is used in the momentum equation; the raw length $L$ is
used for volume accounting.

% =============================================================================
% 3. SIMULATION ORCHESTRATION
% =============================================================================
\section{Simulation orchestration}

\subsection{The overall time-stepping loop}

The engine is a single-step API. Each call to \texttt{SWMMEngine::step()}
executes exactly one routing step (or one runoff-only step when routing is
disabled). The sequence within a step is (\texttt{SWMMEngine.cpp:941}):
\begin{enumerate}[nosep]
  \item Compute the routing time step $dt_{\mathrm{next}}$. If the variable-step
        option is enabled, the Courant-limited step is
        \[
          dt_{\mathrm{cfl}} = \texttt{router.getAdaptiveStep}(
            \texttt{routing\_step},\ \texttt{variable\_step}),
        \]
        then $dt_{\mathrm{next}} = \min(\texttt{routing\_step},\ dt_{\mathrm{cfl}})$
        clamped to the remaining simulation duration.
  \item Save the old hydraulic state (\texttt{ctx.save\_state()}).
  \item Reset the per-step mass-balance accumulators.
  \item Run the pipeline:
        \begin{enumerate}[nosep]
          \item \texttt{stepRunoff($dt$)} --- hydrology (always);
          \item if routing is enabled: \texttt{stepRouting($dt$)} (hydraulics +
                quality), \texttt{updateStatistics($dt$)},
                \texttt{updateRoutingMassBalance($dt$)};
          \item \texttt{computeFinalStorage()}.
        \end{enumerate}
  \item Advance the simulation clock and post output snapshots.
\end{enumerate}

\subsection{Runoff and routing clocks}

The hydrology and hydraulics run on \emph{independent} clocks:
\begin{itemize}[nosep]
  \item The \textbf{runoff clock} advances at the wet step (300\,s default) or
        dry step (3600\,s default), shortened to align with rain-gage
        boundaries. \texttt{stepRunoff} runs multiple runoff sub-steps within
        one routing step ($\text{while } t_{\mathrm{runoff}} <
        t_{\mathrm{routing}} + dt$).
  \item The \textbf{routing clock} advances by $dt$ (the routing step) each
        \texttt{step()} call.
\end{itemize}
After the runoff sub-steps, the wet-weather lateral inflow delivered to nodes
is obtained by \emph{linear interpolation} between the bracketing runoff
evaluations: with $f = (t_{\mathrm{elapsed}} - t_{\mathrm{runoff,old}}) /
(t_{\mathrm{runoff,new}} - t_{\mathrm{runoff,old}})$,
\begin{equation}
  Q_{\mathrm{runoff}} = (1-f)\,Q_{\mathrm{runoff}}^{\mathrm{old}} +
  f\,Q_{\mathrm{runoff}}^{\mathrm{new}},
\end{equation}
plus runon and groundwater contributions. These are scattered into
\texttt{nodes.runoff\_inflow} / \texttt{nodes.gw\_inflow}.

\subsection{Lateral inflows}

Every routing step, the lateral inflows are assembled from decomposed source
buffers in legacy order (\texttt{assembleLateralInflows},
\texttt{SWMMEngine.cpp:5691}):
\begin{equation}
  Q_{\mathrm{lat}} = Q_{\mathrm{ext}} + Q_{\mathrm{dwf}} +
  Q_{\mathrm{wet}} + Q_{\mathrm{gw}} + Q_{\mathrm{rdii}} +
  Q_{\mathrm{iface}} + Q_{\mathrm{user}} + Q_{\mathrm{coupling}},
\end{equation}
where the terms are external/inflow time series, dry-weather flow,
wet-weather (runoff) inflow, groundwater inflow, RDII, interface-file
inflow, user-forced inflow, and 1D--2D coupling inflow respectively.

\subsection{The routing step}

\texttt{Router::step(ctx, dt, evap\_rate, non\_conduit\_fn)}
(\texttt{Routing.cpp:308}) performs, in order:
\begin{enumerate}[nosep]
  \item \textbf{initNodeFlows} --- initialize each node's inflow from lateral
        flow and outflow from losses (storage evaporation and Green--Ampt
        exfiltration, jointly capped against the stored volume); set overflow
        from excess stored volume.
  \item \textbf{computeConduitLosses} --- conduit evaporation (open sections)
        and seepage loss rates. For DYNWAVE this is instead recomputed per
        Picard iteration inside the solver.
  \item \textbf{setAllOutfallDepths} --- set the outfall boundary stages
        (Section~\ref{sec:outfall}).
  \item \textbf{Solver dispatch} --- KINWAVE, DYNWAVE, FV, or STEADY.
  \item \textbf{computeDividerFlows} --- apply flow-divider logic.
  \item \textbf{updateLinkStates} --- recompute node heads $H = z_{\mathrm{inv}}
        + d$.
\end{enumerate}

Inside the dynamic-wave solver, the non-conduit callback computes
pump/orifice/weir/outlet flows inside the Picard iteration, in link-index
order, with immediate scatter into the node inflow/outflow accumulators.

% =============================================================================
% 4. CROSS-SECTION GEOMETRY
% =============================================================================
\section{Cross-section geometry}
\label{sec:xsect}

\subsection{The four geometric relations}

For a conduit, the flow area $A$, top width $W$, wetted perimeter $P_w$, and
hydraulic radius $R = A/P_w$ are functions of the flow depth $y$. The engine
pre-computes, for each shape, the four fundamental relations
(\texttt{XSectKernels.hpp}, \texttt{xsect\_tables.hpp}):
\begin{align}
  A &= A(y), \qquad W = W(y), \qquad R = R(y), \\
  S &= S(A) = A\,R^{2/3} \quad \text{(section factor, ``conveyance shape''),}
\end{align}
together with their inverses $y = y(A)$, $A = A(S)$, and the derivative
$dS/dA$. Closed-form expressions exist for the standard shapes (circular,
filled circular, rectangular, trapezoidal, triangular, parabolic, power,
eggs, horseshoe, gothic, basket-handle, etc.); irregular and custom shapes use
tabulated (transect) tables with linear interpolation, and the inverse
$A(S)$ is solved by Newton iteration or table inversion.

\subsection{Manning and the section factor}

Manning's equation in US customary form is
\begin{equation}
  Q = \frac{1.486}{n}\,A\,R^{2/3}\,\sqrt{S_f}
  = \beta\,\Psi(A)\,\sqrt{S_f/S_0},
  \qquad \Psi(A) = A\,R^{2/3},
\end{equation}
where $S_f$ is the friction slope, $\beta = (1.486/n)\sqrt{S_0}$, and
$\Psi(A)$ is the section factor. For many closed shapes the section factor
peaks below full depth, so the maximum discharge occurs at a less-than-full
area.

% =============================================================================
% 5. STEADY FLOW ROUTING
% =============================================================================
\section{Steady flow routing}

Under steady flow (\texttt{executeSteadyFlow}, \texttt{Routing.cpp:637}), links
are processed in topologically sorted order. For each conduit:
\begin{enumerate}[nosep]
  \item The upstream node's inflow $Q_{\mathrm{in}}$ is gathered (limited by
        the maximum outflow).
  \item The per-barrel flow is $q = Q_{\mathrm{in}}/\mathrm{barrels}$ minus
        the conduit loss rate, clamped at zero.
  \item If $q \ge Q_{\mathrm{full}}$: $q = Q_{\mathrm{full}}$ and
        $A = A_{\mathrm{full}}$; otherwise the normal-depth area is found from
        the inverse section factor,
        \begin{equation}
          s = \frac{q}{\beta}, \qquad A = A(s),
        \end{equation}
        and the depth from $y = y(A)$.
  \item The flow is uniform along the conduit ($Q_{\mathrm{out}} =
        q\cdot\mathrm{barrels}$, $A_1 = A_2 = A$); the link volume is
        $V = A\,L\,\mathrm{barrels}$.
\end{enumerate}
Non-conduit links simply pass the upstream node's inflow through unchanged.
Steady flow converges in one pass.

% =============================================================================
% 6. KINEMATIC WAVE
% =============================================================================
\section{Kinematic wave routing}

\subsection{Governing equations}

The kinematic wave model combines the continuity equation with the uniform-flow
rating curve, dropping the inertia, pressure-gradient, and (partially)
convective terms of the St.\ Venant momentum equation. It cannot represent
backwater, flow reversal, or pressurization, and requires a directed acyclic
network. Starting from the St.\ Venant equations and setting the bed slope
equal to the friction slope, $S_0 = S_f$, the governing pair is:
\begin{equation}
  \frac{\partial A}{\partial t} + \frac{\partial Q}{\partial x} = 0
  \qquad\text{(continuity),}
\end{equation}
\begin{equation}
  Q = \beta\,\Psi(A), \qquad \Psi(A) = A\,R^{2/3},
  \qquad \beta = \frac{1.486}{n}\sqrt{S_0}
  \qquad\text{(rating curve).}
\end{equation}

\subsection{Finite-difference scheme}

A weighted Wendroff implicit finite-difference scheme discretizes the
continuity equation between the upstream (1) and downstream (2) ends:
\begin{equation}
  \frac{(1-\theta)(A_1^{t+1}-A_1^{t}) +
        \theta\,(A_2^{t+1}-A_2^{t})}{\Delta t}
  \;+\;
  \frac{(1-\varphi)(Q_2^{t}-Q_1^{t}) +
        \varphi\,(Q_2^{t+1}-Q_1^{t+1})}{L}
  = 0,
\end{equation}
with $\theta = \varphi = 0.6$. Because each junction has at most one outlet
conduit, processing links in topological order leaves only $A_2^{t+1}$ and
$Q_2^{t+1}$ as unknowns; $Q_1^{t+1}$ is known from the upstream node's inflow
and $A_1^{t+1}$ from the inverse section factor at $Q_1^{t+1}/\beta$.
Substituting the rating curve yields the single nonlinear equation
\begin{equation}
  f\!\left(A_2^{t+1}\right) = \beta\,\Psi\!\left(A_2^{t+1}\right) +
  C_1\,A_2^{t+1} + C_2 = 0,
\end{equation}
with (in the engine's normalized implementation,
\texttt{KinematicWave.cpp:solveConduit})
\begin{align}
  C_1 &= \frac{L\,\theta}{\Delta t\,\varphi},\\
  C_2 &= \frac{L}{\Delta t\,\varphi}
        \Big[(1-\theta)(A_1^{t+1} - A_1^t) - \theta\,A_2^t\Big]
        + \frac{1-\varphi}{\varphi}(Q_2^t - Q_1^t) - Q_1^{t+1}.
\end{align}

\subsection{Root finding}

Equation $f(A_2^{t+1}) = 0$ is solved by a bracketed Newton--Raphson iteration
(maximum 40 iterations). Because the section factor can have two roots
(peak at $A_{\max}$ below full), the bracket is pre-screened: the initial
bracket is $[A_{\max}, A_{\mathrm{full}}]$; if both bounds produce the same
sign it is reset to $[0, A_{\max}]$. If both bounds give negative $f$ the
conduit is full ($Q = Q_{\mathrm{full}}$); if both give positive $f$ the flow
is zero.

\subsection{Storage nodes under kinematic wave}

Storage nodes may have any number of outlet links. Their mass balance is
integrated with the trapezoidal rule,
\begin{equation}
  V^{t+1} = V^{t} + \tfrac{1}{2}\left(Q_{\mathrm{in}}^{t} + Q_{\mathrm{in}}^{t+1}
  - Q_{\mathrm{out}}^{t} - Q_{\mathrm{out}}^{t+1}\right)\Delta t,
\end{equation}
solved by successive approximation ($\omega = 0.55$, convergence tolerance
0.005 ft) because both $Q_{\mathrm{out}}$ and $V$ depend on the head. The
storage head is then updated once more after all link flows are known
(\texttt{Routing.cpp:345--389}).

% =============================================================================
% 7. DYNAMIC WAVE: GOVERNING EQUATIONS
% =============================================================================
\section{Dynamic wave routing: governing equations}

\subsection{The St.\ Venant equations}

The dynamic wave model solves the full one-dimensional de~Saint-Venant
equations (continuity and momentum):
\begin{equation}
  \underbrace{\frac{\partial A}{\partial t}}_{\text{storage}} +
  \underbrace{\frac{\partial Q}{\partial x}}_{\text{convective flux}}
  = 0,
  \label{eq:sv-cont}
\end{equation}
\begin{equation}
  \underbrace{\frac{\partial Q}{\partial t}}_{\text{local acceleration}} +
  \underbrace{\frac{\partial}{\partial x}\!\left(\frac{Q^{2}}{A}\right)}_{\text{convective acceleration}}
  + \underbrace{gA\,\frac{\partial H}{\partial x}}_{\text{pressure gradient}}
  + \underbrace{gA\,S_f}_{\text{friction}}
  = 0,
  \label{eq:sv-mom}
\end{equation}
where $A$ is the flow area, $Q$ the discharge, $H = Z + y$ the hydraulic head
($Z$ conduit invert, $y$ depth), $g$ gravity, and $S_f$ the friction slope.
The friction slope follows Manning:
\begin{equation}
  S_f = \left(\frac{n}{1.486}\right)^{2}
  \frac{Q\,\lvert U \rvert}{A\,R^{4/3}},
  \qquad U = \frac{Q}{A}.
\end{equation}

\subsection{Finite-difference momentum equation}

The engine uses an implicit (backwards Euler) finite-difference form. With
spatial differences over the conduit length $L$ and temporal differences over
$\Delta t$, the flow update at each conduit is
\begin{equation}
  Q^{t+\Delta t}
  = \frac{Q^{t} + \Delta Q_{\mathrm{inertia}} + \Delta Q_{\mathrm{pressure}}
        + \Delta Q_{\mathrm{loss}}}
         {1 + \Delta Q_{\mathrm{friction}}},
  \label{eq:mom-update}
\end{equation}
where all geometric quantities are evaluated at the new time $t + \Delta t$.
The four terms are:
\begin{align}
  \Delta Q_{\mathrm{inertia}}
    &= 2\bar{U}\left(\bar{A}^{t+\Delta t} - \bar{A}^{t}\right)
       + \bar{U}^{2}\,\frac{A_2 - A_1}{L}\,\Delta t, \label{eq:mom-inertia}\\
  \Delta Q_{\mathrm{pressure}}
    &= -g\bar{A}\,\frac{H_2 - H_1}{L}\,\Delta t, \label{eq:mom-pressure}\\
  \Delta Q_{\mathrm{friction}}
    &= g\left(\frac{n}{1.486}\right)^{2}
       \frac{\lvert \bar{U} \rvert}{\bar{R}^{4/3}}\,\Delta t, \label{eq:mom-friction}\\
  \Delta Q_{\mathrm{loss}}
    &= \text{local (minor) losses and evap/seepage terms}.
\end{align}

\subsection{Node continuity}

Conservation of volume at each node assembly (the node plus half of each
connected link) requires
\begin{equation}
  \frac{\partial V}{\partial t}
  = A_{S}\,\frac{\partial H}{\partial t} = \sum Q,
\end{equation}
with $A_S = A_{SN} + \sum A_{SL}$ the assembly surface area (the node's own
storage surface area plus the half-link surface areas contributed by each
conduit). In finite-difference form:
\begin{equation}
  H^{t+\Delta t} = H^{t} +
  \frac{\tfrac{\Delta t}{2}\left(\sum Q^{t} + \sum Q^{t+\Delta t}\right)}
       {\left(A_{SN} + \sum A_{SL}\right)^{t+\Delta t}},
\end{equation}
i.e.\ a trapezoidal update of the head using the \emph{average} of the net
flows at $t$ and $t+\Delta t$.

% =============================================================================
% 8. THE DYNAMIC WAVE SOLUTION ALGORITHM
% =============================================================================
\section{The dynamic wave solution algorithm}

\subsection{Overview: the Picard (successive-approximation) iteration}

Equations~\eqref{eq:mom-update} and the node continuity are solved implicitly
over each time step by \emph{functional iteration} (Picard iteration /
successive approximations). The steps, matching the engine's
\texttt{DWSolver::execute} (\texttt{DynamicWave.cpp:1026}), are:

\begin{enumerate}[nosep]
  \item \textbf{initNodeStates} --- reset per-node accumulators: inflow = 0,
        outflow = losses $+$ negative lateral flow, inflow $+$ positive lateral
        flow; surface area from the storage/ponding geometry; converged = 0,
        $\sum dQ/dH = 0$.
  \item \textbf{computeLinkGeometry} --- batch-compute the end depths
        $y_1, y_2$, midpoint depth $\bar y$, top widths, areas, and hydraulic
        radii of all conduits from the current node heads; classify each
        conduit's flow regime; apply Preissmann-slot overrides where
        surcharged.
  \item \textbf{Momentum kernels} --- solve the implicit momentum equation for
        each conduit (Section~\ref{sec:momentum-kernel}), commit the new
        flows, and scatter them to the node inflow/outflow accumulators,
        accumulating $\sum dQ/dH$ and the surface-area contributions.
  \item \textbf{Non-conduit structures} --- compute pump/orifice/weir/outlet
        flows (Section~\ref{sec:structures}).
  \item \textbf{Outfall depths} --- re-set outfall boundary stages from the
        current flows.
  \item \textbf{Node depth update} --- solve the continuity equation for each
        node head (Section~\ref{sec:node-continuity}); tally the unconverged
        nodes.
  \item If any node changed by more than the head tolerance and the trial limit
        has not been reached, mark links whose both end nodes converged as
        \emph{bypassed} (their flows are held), and iterate.
\end{enumerate}

Under-relaxation: after the first iteration every updated flow and head is
relaxed with $\omega = 0.5$,
\begin{equation}
  x^{k+1} = (1-\omega)\,x^{k} + \omega\,x^{k+1}_{\mathrm{raw}},
  \qquad \omega = 0.5,
\end{equation}
and a sign change in flow is forced through a small nonzero value
($q = \pm 0.001$).

\subsection{Convergence}

A node is converged when both the raw Picard residual and the accepted (possibly
Anderson-mixed) movement are within the head tolerance $\varepsilon_H = 0.005\
\ft$ (default). The routing step converges when no non-outfall node is
unconverged. The maximum number of trials is 8 by default
(\texttt{MAX\_TRIALS}); a step that converges on its last allowed trial is still
counted as converged.

\subsection{Flow classification}

During geometry computation each conduit is classified into one of seven flow
classes (DRY, UP\_DRY, DN\_DRY, SUBCRITICAL, SUPERCRITICAL, UP\_CRITICAL,
DN\_CRITICAL) based on the end depths and the head relative to the end invert
elevations, using the normal depth $y_n$ (from the inverse section factor at
$Q/\beta$) and critical depth $y_c$ (Section~\ref{sec:outfall}). The class
drives:
\begin{itemize}[nosep]
  \item the surface-area adjustments for dry/critical ends (Table~\ref{tab:surfadj}),
  \item the depth used at the critical end (replaced by $\min(y_n, y_c)$),
  \item the friction-based normal-flow limit.
\end{itemize}

\begin{table}[H]
\centering
\small
\begin{tabular}{lll}
\toprule
\textbf{Condition} & \textbf{Criteria} & \textbf{Adjustment} \\
\midrule
Upstream dry & $y_1 = 0$; $Z_1 > E_1$ & $A_{SL1} = 0$ if $H_2 \le Z_1$, else critical adjustment \\
Downstream dry & $y_2 = 0$; $Z_2 > E_2$ & $A_{SL2} = 0$ if $H_1 \le Z_2$, else critical adjustment \\
Upstream critical & $Q < 0$; $Z_1 > E_1$; $H_1 - Z_1 < y_*$ & $y_1 = y_*$; $A_{SL1} = 0$ \\
Downstream critical & $Q > 0$; $Z_2 > E_2$; $H_2 - Z_2 < y_*$ & $y_2 = y_*$; $A_{SL2} = 0$ \\
\bottomrule
\end{tabular}
\caption{Surface-area and depth adjustments for dry/critical conduit ends
($y_* = \min(y_n, y_c)$; $E$ = node invert, $Z$ = conduit invert).}
\label{tab:surfadj}
\end{table}

\subsection{Surface-area contributions}

For a subcritical conduit, the free-surface areas contributed to its two end
nodes are
\begin{equation}
  A_{SL1} = \frac{W_1 + \bar{W}}{2}\,\frac{L}{2},
  \qquad
  A_{SL2} = \frac{\bar{W} + W_2}{2}\,\frac{L}{2}\,
  \mathrm{fasnh},
\end{equation}
where $W_1, W_2, \bar W$ are the top widths at $y_1, y_2, \bar y$, and
$\mathrm{fasnh}$ is a factor between normal and critical depth for a nearly
critical downstream end. For dry/critical ends the half-link surface area is
adjusted per Table~\ref{tab:surfadj}; a completely dry conduit contributes
$A_{SL1} = A_{SL2} = \mathrm{FUDGE}\cdot L/2$. A closed conduit whose midpoint
depth exceeds the crown cutoff ($y/y_{\mathrm{full}} \ge 0.96$ under EXTRAN)
uses a width capped at the crown in the surface-area computation, providing the
small nonzero contribution that keeps the node Picard denominator bounded.

% =============================================================================
% 9. THE MOMENTUM KERNEL
% =============================================================================
\section{The link momentum kernel}
\label{sec:momentum-kernel}

Each conduit is solved by one of the momentum kernels in
\texttt{DynamicWave.cpp} (\texttt{process\-Manning\-Link},
\texttt{process\-Dry\-Link}, \texttt{process\-Force\-Main\-Link}), dispatched
by momentum category. This section transcribes the Manning (free-surface and
closed-conduit) kernel, which is the main one.

\subsection{Velocity, Froude number, and inertial damping}

With $\bar A$ the midpoint area and $q_{\mathrm{last}}$ the per-barrel flow
from the previous iteration,
\begin{equation}
  v = \frac{q_{\mathrm{last}}}{\bar A}, \qquad
  \lvert v \rvert \le V_{\max} = 50\ \ft/\st
  \quad\text{(velocity cap).}
\end{equation}
For a non-surcharged conduit the Froude number uses the hydraulic depth
$\bar A/\bar W$:
\begin{equation}
  \mathrm{Fr} = \frac{\lvert v \rvert}{\sqrt{g\,\bar A / \bar W}}
  \quad\text{(0 for a closed conduit within FUDGE of full).}
\end{equation}
The inertial-damping factor follows the \emph{local partial inertia} approach
(linear blend over $0.5 \le \mathrm{Fr} \le 1$):
\begin{equation}
  \sigma =
  \begin{cases}
    1 & \mathrm{Fr} \le 0.5,\\
    2\,(1 - \mathrm{Fr}) & 0.5 < \mathrm{Fr} < 1,\\
    0 & \mathrm{Fr} \ge 1,
  \end{cases}
  \qquad\text{i.e.}\quad
  \sigma = \mathrm{clamp}\!\big(2(1-\mathrm{Fr}),\ 0,\ 1\big).
\end{equation}
The \texttt{INERTIAL\_DAMPING} option overrides this: NONE forces $\sigma = 1$,
FULL forces $\sigma = 0$ (no inertial terms at all); a closed conduit that is
surcharged always has $\sigma = 0$.

\subsection{Upstream weighting}

The pressure and friction terms use \emph{upstream-weighted} average area and
hydraulic radius, reflecting that supercritical flow is influenced only by
upstream conditions. With the weighting factor
\begin{equation}
  \rho = \begin{cases} \sigma & q_{\mathrm{last}} > 0 \text{ and } H_1 \ge H_2
    \text{ (and not full),}\\ 1 & \text{otherwise,}\end{cases}
\end{equation}
\begin{equation}
  \bar A' = A_1 + \rho\left(\bar A - A_1\right),
  \qquad
  \bar R' = R_1 + \rho\left(\bar R - R_1\right).
\end{equation}

\subsection{The six momentum terms}

The engine evaluates the momentum equation as the six contributions
$dq_1 \ldots dq_6$ (per barrel):
\begin{align}
  dq_1 &= \Delta t \cdot \mathrm{rough\_factor} \cdot
        \frac{\lvert v \rvert}{\bar R'^{\,4/3}}, \label{eq:dq1}\\
  dq_2 &= \Delta t\, g\, \bar A'\, \frac{H_2 - H_1}{L}, \label{eq:dq2}\\
  dq_3 &= 2\,v\,\sigma\left(\bar A^{t+\Delta t} - \bar A^{t}\right), \label{eq:dq3}\\
  dq_4 &= \Delta t\, v^{2}\, \sigma\,\frac{A_2 - A_1}{L}, \label{eq:dq4}\\
  dq_5 &= \frac{\Delta t}{2L}
          \left[ K_{\mathrm{in}}\frac{\lvert q \rvert}{A_1} +
                K_{\mathrm{out}}\frac{\lvert q \rvert}{A_2} +
                K_{\mathrm{avg}}\frac{\lvert q \rvert}{\bar A} \right], \label{eq:dq5}\\
  dq_6 &= \frac{2.5\,\Delta t\, v}{L}\left(Q_{\mathrm{evap}} +
          Q_{\mathrm{seep}}\right), \label{eq:dq6}
\end{align}
where $\mathrm{rough\_factor} = g(n/1.486)^2$, the local-loss coefficients
$K_{\mathrm{in}}, K_{\mathrm{out}}, K_{\mathrm{avg}}$ are the user's inlet,
outlet, and average (loss-weighted) minor-loss coefficients, and
$Q_{\mathrm{evap}} + Q_{\mathrm{seep}}$ is the per-conduit loss rate.

The flow update (per barrel) and its head gradient are
\begin{equation}
  q = \frac{q_{\mathrm{old}} - dq_2 + dq_3 + dq_4 + dq_6}
           {1 + dq_1 + dq_5},
  \label{eq:q-update}
\end{equation}
\begin{equation}
  \frac{dQ}{dH} = \frac{1}{1 + dq_1 + dq_5}\,
  \frac{g\,\Delta t\, \bar A'}{L}\,\mathrm{barrels}.
  \label{eq:dqdh}
\end{equation}
The $dQ/dH$ gradient feeds the node surcharge solver (Section~\ref{sec:node-cont}).

\subsection{Flow limiting and post-processing}

After the raw momentum update, \texttt{applyFlowLimits}
(\texttt{DynamicWave.cpp:2211}) applies, in order:
\begin{enumerate}[nosep]
  \item \textbf{Culvert inlet control} (FHWA HEC-5): if a culvert code is
        present and the conduit is not full, $q \leftarrow \min(q, q_{\mathrm{inlet}})$.
  \item \textbf{Normal-flow limit}: for an open/free-surface conduit not full,
        if the slope condition ($y_1 < y_2$, i.e.\ the water-surface slope is
        smaller than the bed slope) or the upstream-Froude condition
        ($\mathrm{Fr}_1 \ge 1$) holds (per the \texttt{NORMAL\_FLOW\_LIMITED}
        option), then
        \begin{equation}
          q \leftarrow \min\!\left(q,\ \beta\,A_1\,R_1^{2/3}\right),
        \end{equation}
        the Manning normal flow at the upstream depth.
  \item \textbf{Under-relaxation} (iterations $> 0$):
        $q = (1-\omega)q_{\mathrm{last}} + \omega\,q$ with a sign-change clamp
        to $\pm 0.001$.
  \item \textbf{User flow limit}: $\lvert q \rvert \le q_{\mathrm{limit}}$.
  \item \textbf{Flap gates}: reverse flow through a flap gate is zeroed.
  \item \textbf{Dry-node check}: flow is forced to $\pm \mathrm{FUDGE}$ if the
        upstream (downstream) node is dry.
\end{enumerate}
Finally,
\begin{equation}
  Q^{t+\Delta t} = q\,\mathrm{barrels}, \qquad
  d_{\mathrm{link}} = \min(\bar y, y_{\mathrm{full}}), \qquad
  V_{\mathrm{link}} = \frac{A_1 + A_2}{2}\,L_{\mathrm{raw}}\,\mathrm{barrels}.
\end{equation}

\subsection{Force mains}

Force mains (a \texttt{FORCE\_MAIN} cross-section, always full) use a separate
kernel (\texttt{process\-Force\-Main\-Link}) with $\sigma = 0$ and friction
from either the Hazen--Williams or Darcy--Weisbach equation,
\begin{equation}
  \begin{aligned}
    dq_1 &= \Delta t\, g\, \frac{S_f}{\lvert v \rvert},\\
    q &= \frac{q_{\mathrm{old}} - dq_2 + dq_6}{1 + dq_1 + dq_5}.
  \end{aligned}
\end{equation}

% =============================================================================
% 10. NODE CONTINUITY AND SURCHARGE
% =============================================================================
\section{Node continuity, flooding, and surcharge}
\label{sec:node-cont}

\subsection{Net flow and volume change}

At each node, per Picard iteration (\texttt{DynamicWave.cpp}, function
\texttt{setNodeDepth}), the net flow and trapezoidal volume change are:
\begin{equation}
  \Delta Q = Q_{\mathrm{in}} - Q_{\mathrm{out}}, \qquad
  \Delta V = \tfrac{1}{2}\left(\Delta Q_{\mathrm{net}}^{t-1} + \Delta Q\right)
  \Delta t,
\end{equation}
where $\Delta Q_{\mathrm{net}}^{t-1}$ is the net flow of the previous step
(trapezoidal/Crank--Nicolson averaging).

\subsection{Free-surface update (EXPLICIT continuity)}

For a node that is not surcharged, the depth update is
\begin{equation}
  \Delta y = \frac{\Delta V}{A_S}, \qquad
  y^{t+\Delta t} = y^{t} + \Delta y,
\end{equation}
where the assembly surface area is floored by the minimum surface area,
\begin{equation}
  A_S = \max\left(A_{SN} + \sum A_{SL},\ A_{\min}\right),
  \qquad A_{\min} = 12.566\ \sfft,
\end{equation}
the default (overrideable via the \texttt{MIN\_SURFAREA} option).
The minimum area is a purely computational device; it does not add volume.

\subsection{Surcharge (EXTRAN) update}

A node is surcharged when all conduits connected to it are full or its head
exceeds the crown of its highest conduit ($y > y_{\mathrm{crown}}$), and it is
not a storage or outfall node and is not ponding. Under EXTRAN the continuity
equation becomes $\sum Q = 0$, enforced through a perturbation (Newton/Hardy
Cross) form. The surcharged depth update is
\begin{equation}
  \Delta y = \frac{\alpha\,\Delta Q}{\mathrm{denom}}, \qquad
  y^{t+\Delta t} = y^{t} + \Delta y,
\end{equation}
with the terminal-node correction $\alpha = 0.6$ for upstream terminal nodes
(only outflow links, $\mathrm{deg} < 0$) and 1 otherwise, and a crown-proximity
blend for a smooth transition into/out of surcharge:
\begin{equation}
  \mathrm{denom} = \sum \frac{dQ}{dH} +
  \left(\frac{A_S^{\mathrm{old}}}{\Delta t} - \sum \frac{dQ}{dH}\right)
  e^{-15\,f},
  \qquad f = \frac{y^{t} - y_{\mathrm{crown}}}{y_{\mathrm{crown}}},
\end{equation}
applied while $y^{t} < 1.25\,y_{\mathrm{crown}}$. The head is kept at or above
the crown ($y \ge y_{\mathrm{crown}} - \mathrm{FUDGE}$). The $\sum dQ/dH$
accumulated at each node is the sum of the per-link gradients
\eqref{eq:dqdh}; note the engine accumulates $\sum dQ/dH$ as a positive
magnitude with $dQ_{\mathrm{net}}/dH = -\sum dQ/dH$ (a head rise increases net
outflow).

\subsection{Semi-implicit node continuity}

The \texttt{NODE\_CONTINUITY} option (default EXPLICIT) can select the unified
\emph{semi-implicit} formulation, a single smooth expression valid for both
free-surface and surcharged states. Linearizing the net flow about the current
head estimate with the flow gradients, $\sum Q^{t+\Delta t} \approx \sum Q +
\sum(dQ/dH)\,\Delta H$, and substituting into the trapezoidal head update
yields
\begin{equation}
  \Delta y = \frac{\Delta V}
    {A_S + \tfrac{\Delta t}{2}\sum \frac{dQ}{dH}}, \qquad
  y^{t+\Delta t} = y^{t} + \Delta y,
\end{equation}
with the denominator floored by $A_{\min}$. The flow-gradient term damps the
update (a head rise that increases net outflow reduces the net fill), and the
expression is $C^1$-smooth through the crown, which is what makes it compatible
with Anderson acceleration.

\subsection{Flooding and ponding}

The maximum non-flooded depth is $y_{\max} = d_{\mathrm{full}} + d_{\mathrm{sur}}$
(for non-ponding nodes; ponding nodes are never capped). If the candidate depth
exceeds $y_{\max}$:
\begin{itemize}[nosep]
  \item \textbf{Non-ponding}: the head is capped at $y_{\max}$ and the excess is
        lost as overflow,
        \begin{equation}
          Q_{\mathrm{ovfl}} = \frac{\Delta V}{\Delta t}, \qquad
          V = V_{\mathrm{full}};
        \end{equation}
  \item \textbf{Ponding} (with \texttt{ALLOW\_PONDING} and a positive ponded
        area): the ponded node acts as a constant-area storage node above the
        rim,
        \begin{equation}
          V = \max\left(V^{\mathrm{old}} + \Delta V,\ V_{\mathrm{full}}\right), \qquad
          Q_{\mathrm{ovfl}} = \frac{V - \max(V^{\mathrm{old}},\ V_{\mathrm{full}})}
            {\Delta t},
        \end{equation}
        and the head is free to rise above the rim.
\end{itemize}
A ponded node is not allowed to drop much below full depth once ponded
($y \ge d_{\mathrm{full}} - \mathrm{FUDGE}$).

% =============================================================================
% 11. SURCHARGE METHODS
% =============================================================================
\section{Surcharge methods}

The \texttt{SURCHARGE\_METHOD} option selects how pressurized flow in closed
conduits is treated. All three methods update the flow area of a surcharged
conduit (depth $> y_{\mathrm{full}}$) by adding a narrow slot of width $w_s$:
\begin{equation}
  A = A_{\mathrm{full}} + (y - y_{\mathrm{full}})\,w_s,
\end{equation}
while the hydraulic radius is clamped to its full value $R_{\mathrm{full}}$
(the slot contributes storage but not friction).

\subsection{EXTRAN (default)}

The classical perturbation method of Section~\ref{sec:node-cont}:
$dQ/dH$-based surcharged node updates with a crown cutoff at
$y/y_{\mathrm{full}} = 0.96$ used for the width used in the surface-area
computation.

\subsection{Static Preissmann slot}

The conduit carries a fictitious narrow slot above the crown so the free-surface
formulation remains valid throughout pressurization. The slot width follows the
Sj{\"o}berg (1982) formula,
\begin{equation}
  \frac{w_s}{W_{\max}} = 0.5423\,
  \exp\!\left(-\left(\frac{y}{y_{\mathrm{full}}}\right)^{2.4}\right),
\end{equation}
applied for $0.985257 \le y/y_{\mathrm{full}} \le 1.78$, clamped at
$w_s = 0.01\,W_{\max}$ above that, and zero below the cutoff.

\subsection{Dynamic Preissmann slot}

An OpenSWMM extension (Sharior, Hodges \& Vasconcelos 2023) in which the slot
area evolves in time as an element of transient storage. The slot top width is
driven by a target pressure-wave celerity $c_{pT}$ and a time-varying
Preissmann number $P$:
\begin{equation}
  T_s = \frac{g\,A_{\mathrm{full}}}{c_{pT}^{2}}\,P^{2},
  \qquad
  P_{0} = \max\!\left(\frac{c_{pT}}{\alpha\,c_g},\ 1\right),
  \qquad c_g = \sqrt{g\,A_{\mathrm{full}}/W_{\max}},
\end{equation}
with the shock parameter $\alpha = 3$ (default). The accumulated slot area is
path-dependent,
\begin{equation}
  A_s \leftarrow \max\!\left(A_s + T_s\,\Delta h_s,\ 0\right),
  \qquad h_s = \max(\bar y - y_{\mathrm{full}},\ 0),
\end{equation}
and $P$ decays exponentially after pressurization onset and is spatially
smoothed across nodes. While a slot is active, the top width is $T_s$ and the
surface area contributed to a surcharged end node is $T_s\,L/4$. Nodal heads
are updated with the ordinary free-surface formula at all times; the
surcharge branch is never invoked. The variable time step uses the pressure
celerity $c_p = c_{pT}/P$.

% =============================================================================
% 12. NON-CONDUIT STRUCTURES
% =============================================================================
\section{Non-conduit hydraulic structures}
\label{sec:structures}

The non-conduit structure flows are computed by \texttt{HydStructures.cpp}
(\texttt{StructureSolver}) inside the dynamic-wave Picard iteration, one link
at a time in link-index order (so that two pumps sharing a wet well see each
other's draw), with immediate scatter into the node accumulators and
under-relaxation at $\omega = 0.5$ for iterations $> 0$ (pumps exempt).

\subsection{Pumps}

Pump on/off control uses depth hysteresis applied once per routing step:
\begin{equation}
  \text{shut off if } d < d_{\mathrm{off}} \ (\text{setting} > 0),
  \qquad
  \text{start if } d > d_{\mathrm{on}} \ (\text{setting} = 0).
\end{equation}
The pump curve types (1--5) determine the discharge $q$ from the curve versus
wet-well volume (type 1), depth (types 2 and 4), or head (types 3 and 5, where
$q$ is scaled by the speed setting $s$ and the head by $s^2$); an ideal pump
discharges exactly the upstream node's inflow plus overflow. Flow limiting
(legacy \texttt{getModPumpFlow}) prevents pumping more than is available
(dry-draw protection) and, for type-1 pumps or storage-fed pumps, caps $q$ by
$Q_{\mathrm{in}} + V_{\mathrm{old}}/\Delta t$.

\subsection{Orifices}

Bottom and side orifices use a coefficient and a critical depth set up from the
opening height $h_{\mathrm{open}} = s\,y_{\mathrm{full}}$ (setting $s$):
\begin{align}
  f &= \min\!\left(\frac{\text{head}}{h_{\mathrm{crit}}},\ 1\right),
  \qquad h_{\mathrm{crit}} = \frac{C_d}{0.414}\,\frac{h_{\mathrm{open}}}
    {4}\ \text{(circular)}\ \text{or}\ \frac{C_d}{0.414}\,
    \frac{h_{\mathrm{open}}\,W_{\max}}{2(h_{\mathrm{open}} + W_{\max})}
    \text{(rectangular)},\\
  q &= \begin{cases}
    C_w\,f^{3/2} & f < 1 \quad\text{(weir-like partial flow)},\\
    C_d\,A_{\mathrm{eff}}\sqrt{2g\,H} & f \ge 1 \quad\text{(full orifice)},
  \end{cases}
\end{align}
with $C_w = C_d\sqrt{h_{\mathrm{crit}}}\,A_{\mathrm{eff}}\sqrt{2g}$ (sharp-
crested weir coefficient 0.414), $A_{\mathrm{eff}}$ the cross-section area at
the opening height, and $H$ the appropriate head (differential head when
submerged). The flow gradient is
\begin{equation}
  \frac{dQ}{dH} = \frac{3}{2}\,\frac{q}{f\,h_{\mathrm{crit}}}
  \quad\text{(weir regime)}, \qquad
  \frac{dQ}{dH} = \frac{q}{2H}
  \quad\text{(orifice regime)}.
\end{equation}
An optional ARMCO flap-gate head loss and the Villemonte submergence correction
may be applied.

\subsection{Weirs}

The four weir types (transverse, side-flow, V-notch, trapezoidal) discharge
according to
\begin{equation}
  \begin{aligned}
    q &= C_d\,L\,h^{3/2} &&\text{(transverse)},\\
    q &= C_d\,L^{0.83}\,h^{1.67} &&\text{(side-flow, corrected)},\\
    q &= C_d\,s_h\,h^{5/2} &&\text{(V-notch)},
  \end{aligned}
\end{equation}
with end contractions reducing the effective length. When the upstream
hydraulic grade line reaches the crown ($H_1 \ge H_{\mathrm{crown}}$) and
surcharging is allowed, the weir transitions to orifice flow using an
equivalent-orifice coefficient computed from the weir discharge at full
opening:
\begin{equation}
  C_{\mathrm{sur}} = \frac{Q_{\mathrm{weir}}(s\,y_{\mathrm{full}})}
    {\sqrt{(s\,y_{\mathrm{full}})/2}},
  \qquad
  q = C_{\mathrm{sur}}\sqrt{H_{\mathrm{orif}}},
\end{equation}
with $H_{\mathrm{orif}}$ the head to the mid-point of the opening (or the
differential head when submerged). If surcharging is not allowed, the head is
capped at the crown and the weir equation is kept. Villemonte submergence
correction applies when the downstream HGL exceeds the crest. The weir's
surface-area contribution is zeroed for SWMM-4 compatibility.

\subsection{Outlets}

Outlets use head- or depth-based rating curves (tabular or functional):
\begin{equation}
  q = C\,H^{e} \quad\text{(functional)}, \qquad
  q = \mathrm{table}(H) \quad\text{(tabular)},
\end{equation}
scaled by the setting. The $dQ/dH$ is deliberately left zero (matching
legacy), which keeps the node surcharge denominator from being inflated.

% =============================================================================
% 13. OUTFALL BOUNDARY CONDITIONS
% =============================================================================
\section{Outfall boundary conditions}
\label{sec:outfall}

Outfall stages are set at the start of each routing step and again inside each
Picard iteration from the current conduit flows. For the conduit connected to
an outfall at end offset $z$, with per-barrel flow $q$, the normal depth
$y_n$ (from the inverse section factor) and critical depth $y_c$ are computed
($y_c$ has closed forms for standard shapes and a numerical root find
otherwise). The stage depends on the boundary-condition type
(\texttt{Outfall.cpp:setOutfallDepth}):
\begin{itemize}[nosep]
  \item \textbf{FREE}: free overfall, $d = z + \min(y_n, y_c)$.
  \item \textbf{NORMAL}: $d = z + y_n$.
  \item \textbf{FIXED}, \textbf{TIDAL}, \textbf{TIMESERIES}: the stage
        $s_{\mathrm{stage}}$ is constant, from a tidal table (hour of day), or
        from a time series; then
        \begin{equation}
          d = \begin{cases}
            s_{\mathrm{stage}} - z_{\mathrm{inv}} & z + y_c + z_{\mathrm{inv}}
              < s_{\mathrm{stage}},\\
            \max(0,\ s_{\mathrm{stage}} - z_{\mathrm{inv}}) & z > 0 \wedge
              s_{\mathrm{stage}} < z_{\mathrm{inv}} + z,\\
            z + y_c & z > 0 \wedge s_{\mathrm{stage}} \ge z_{\mathrm{inv}} + z,\\
            y_c & z = 0.
          \end{cases}
        \end{equation}
\end{itemize}
The last case keeps a free-discharging fixed outfall at the conduit critical
depth. A flap gate at an outfall blocks reverse flow.

% =============================================================================
% 14. FLOW DIVIDERS
% =============================================================================
\section{Flow dividers}

Dividers are active only under kinematic-wave (and steady) routing; under
dynamic wave they act as ordinary junctions. Given the total inflow
$Q_{\mathrm{in}}$ at the divider node:
\begin{itemize}[nosep]
  \item \textbf{Cutoff}: diverts all inflow above a minimum,
        $Q_{\mathrm{div}} = \max(0, Q_{\mathrm{in}} - q_{\min})$.
  \item \textbf{Overflow}: diverts the inflow above the capacity of the
        continuation link.
  \item \textbf{Tabular}: $Q_{\mathrm{div}} = Q_{\mathrm{in}}\cdot
        f(Q_{\mathrm{in}})$ with the fraction from a rating table.
  \item \textbf{Weir}: $Q_{\mathrm{div}} = C_d\,W\,d^{3/2}$ from a weir relation
        on the node depth, capped at the inflow.
\end{itemize}
The diverted flow is written to the diversion link; the continuation link
carries the remainder.

% =============================================================================
% 15. ADAPTIVE TIME STEPPING
% =============================================================================
\section{Adaptive time stepping and the Courant condition}

\subsection{Courant--Friedrichs--Lewy condition}

Because the dynamic-wave scheme updates flows and heads element by element
(there is no simultaneous spatial coupling), it is subject to the Courant
condition: the time step must not exceed the time for a dynamic wave to
traverse the shortest conduit,
\begin{equation}
  \Delta t \le \frac{L}{\lvert U \rvert + c},
  \qquad c = \sqrt{g\,\frac{A}{W}}
  \text{ (gravity-wave celerity).}
\end{equation}

\subsection{Link-based step}

For each conduit (\texttt{getLinkStep}, \texttt{DynamicWave.cpp:3711}) with
per-barrel flow $q = \lvert Q \rvert/\mathrm{barrels}$, midpoint area $A$, and
Froude number $\mathrm{Fr}$:
\begin{equation}
  t = \frac{V/\mathrm{barrels}}{q}\,
      \frac{L'}{L}\,\frac{\mathrm{Fr}}{1 + \mathrm{Fr}},
\end{equation}
(the classic SWMM CFL expression, scaled by the lengthened/raw length ratio).
Under the dynamic Preissmann slot a surcharged conduit instead uses the
pressure celerity,
\begin{equation}
  t = \frac{L\,(L'/L)}{\lvert v \rvert + c_p}, \qquad
  c_p = \frac{c_{pT}}{P}.
\end{equation}

\subsection{Node-based step}

For each non-outfall node below the crown with a nonzero depth-change rate
$\dot y = \lvert \Delta y/\Delta t \rvert$:
\begin{equation}
  t = 0.25\,\frac{y_{\mathrm{crown}}}{\dot y},
\end{equation}
guarding against excessive head change in one step.

\subsection{Combining the constraints}

The next routing step (\texttt{getRoutingStep}) is the minimum over all links,
nodes, and virtual-junction pairs, each scaled by the user's Courant factor,
then floored and clamped:
\begin{equation}
  \Delta t_{\min}^{\mathrm{eff}} =
    \max\!\left(\min_{j,i}\{t_j, t_i\}\cdot \mathrm{Cr},\
      \max(\min(\Delta t_{\min},\ \Delta t_{\mathrm{routing}}),\
      0.001\ \st)\right),
\end{equation}
rounded to milliseconds. The effective step is then capped by the user's fixed
\texttt{ROUTING\_STEP} and the remaining simulation duration. The initial step
is the minimum step. The first call (no flows yet) also returns the minimum
step. The FV solver substeps internally at its own CFL limit, so the routing
step under FV is only a reporting cadence.

% =============================================================================
% 16. STABILITY, CONVERGENCE, MASS BALANCE
% =============================================================================
\section{Stability, convergence, and mass balance}

\subsection{Stability indicators}

Numerical instability appears as non-damping oscillations in flow and water
surface, and nodes that repeatedly dry up. Two report metrics quantify it:
\begin{itemize}[nosep]
  \item the overall \textbf{flow continuity error} (below), and
  \item the \textbf{link Flow Instability Index (FII)} --- the normalized count
        of times a link's flow exceeds both its neighbors.
\end{itemize}

\subsection{Convergence of the Picard loop}

A routing step ``does not converge'' only when the head tolerance
($\varepsilon_H = 0.005\ \ft$) is not met at some non-outfall node after
\texttt{MAX\_TRIALS} iterations. Anderson acceleration (optional) speeds
convergence by blending the two most recent operator outputs at each node:
\begin{equation}
  \alpha_k = \mathrm{clamp}\!\left(
    \frac{r_k\,(r_k - r_{k-1})}{(r_k - r_{k-1})^{2}},\ 0,\ 1\right),
  \qquad
  H_{k+1} = (1-\alpha_k)\,G(H_k) + \alpha_k\,G(H_{k-1}),
\end{equation}
with $r_k = G(H_k) - H_k$ the residual and $G$ the head-update operator. It is
skipped at nodes where $G$ is non-smooth (EXTRAN-surcharged, dynamic-slot
active, near the static-slot cutoff, weir/orifice at crown, pump ends, ponding
boundary).

\subsection{Routing mass balance}

The routing continuity error (fraction of total inflow) is
\begin{equation}
  \varepsilon_{\mathrm{routing}} = \frac{Q_{\mathrm{dw}} + Q_{\mathrm{wet}} +
    Q_{\mathrm{gw}} + Q_{\mathrm{rdii}} + Q_{\mathrm{ext}} +
    V_{\mathrm{init}} - \left(Q_{\mathrm{flood}} + Q_{\mathrm{out}} +
    Q_{\mathrm{evap}} + Q_{\mathrm{seep}} + V_{\mathrm{final}}\right)}
    {Q_{\mathrm{total\,in}}},
\end{equation}
where the inflows are dry-weather, wet-weather, groundwater, RDII, and external,
and the outflows are flooding, system outflow, evaporation, seepage, and final
storage. Flooding is counted only when the node volume does not exceed its full
volume; outfall discharge counts as system outflow, and outfall backflow counts
as external inflow. Conduit and storage-node evap/seep losses are charged here
as well.

% =============================================================================
% 17. OPTIONS
% =============================================================================
\section{Options and defaults}

Table~\ref{tab:options} summarizes the routing-related options
(\texttt{SimulationOptions.hpp}) and their defaults.

\begin{longtable}{lll}
\toprule
\textbf{Option} & \textbf{Default} & \textbf{Purpose} \\
\midrule
\endhead
\texttt{FLOW\_ROUTING} & DYNWAVE & routing formulation \\
\texttt{ROUTING\_STEP} & 20 s & fixed routing step \\
\texttt{MINIMUM\_STEP} & 0.5 s & floor for the CFL step \\
\texttt{DRY\_STEP} / \texttt{WET\_STEP} & 3600 / 300 s & runoff step bounds \\
\texttt{VARIABLE\_STEP} & 0.75 & Courant factor (0 = fixed step) \\
\texttt{LENGTHENING\_STEP} & 0 & Courant lengthening step \\
\texttt{MAX\_TRIALS} & 8 & Picard trial limit \\
\texttt{HEAD\_TOLERANCE} & 0.005 (project units) & node convergence tolerance \\
\texttt{SURCHARGE\_METHOD} & EXTRAN & EXTRAN / SLOT / DYNAMIC\_SLOT \\
\texttt{NODE\_CONTINUITY} & EXPLICIT & two-branch vs semi-implicit \\
\texttt{INERTIAL\_DAMPING} & PARTIAL & Froude damping policy \\
\texttt{NORMAL\_FLOW\_LIMITED} & BOTH & slope/Froude normal-flow cap \\
\texttt{MIN\_SURFAREA} & 12.566 ft$^2$ & nodal area floor \\
\texttt{ALLOW\_PONDING} & false & allow ponding above the rim \\
\texttt{ANDERSON\_ACCEL} & false & Anderson acceleration \\
\texttt{SKIP\_STEADY\_STATE} & false & skip steady routing \\
\bottomrule
\caption{Routing-related options and defaults.}
\label{tab:options}
\end{longtable}

% =============================================================================
% 18. REFERENCES
% =============================================================================
\section{References}

\begin{enumerate}[nosep,label={[\arabic*]}]
  \item Rossman, L.~A. (2017). \emph{Storm Water Management Model Reference
        Manual Volume II -- Hydraulics}. U.S.\ EPA.
  \item OpenSWMM engine source:
        \begin{itemize}[nosep,itemsep=0pt]
          \item \src{hydraulics/Routing.cpp} \\
                \src{hydraulics/DynamicWave.cpp} \\
                \src{hydraulics/KinematicWave.cpp} \\
                \src{hydraulics/HydStructures.cpp}
          \item \src{hydraulics/Node.cpp} \\
                \src{hydraulics/Link.cpp} \\
                \src{hydraulics/Outfall.cpp} \\
                \src{hydraulics/Divider.cpp}
          \item \src{hydraulics/XSectBatch.cpp} \\
                \src{hydraulics/XSectKernels.hpp} \\
                \src{core/SimulationOptions.hpp} \\
                \src{core/SWMMEngine.cpp}
        \end{itemize}
  \item Sharior, S., Hodges, B.~R., \& Vasconcelos, J.~G. (2023).
        Generalized, dynamic, and transient-storage form of the Preissmann
        slot. \emph{Journal of Hydraulic Engineering}, 149(11).
  \item Sj{\"o}berg, A. (1982). Sewer network models dagnum and DIVISION --
        a brief description. \emph{Swedish Water and Wastewater Works
        Association}.
  \item Roesner, L.~A., Nichandros, H.~M., Shubinski, R.~P., Feldman, A.~D.,
        Abbott, J.~W., \& Delleur, J.~W. (1981). \emph{Physical model for
        combined sewer overflows}. Proc.\ ASCE.
  \item de~Almeida, G.~A.~M., \& Bates, P.~D. (2013). Applicability of the
        local inertial approximation of the shallow water equations to flood
        modeling. \emph{Water Resources Research}, 49(8).
\end{enumerate}

\end{document}
